\title[Causality]{Causality: Lecture 13}
\author[Joris Mooij]{\\[5mm]\large Joris Mooij \& Patrick Forr\'e\\
\url{j.m.mooij@uva.nl}, \url{p.d.forre@uva.nl}}
\institute[UvA]{University of Amsterdam}
\date[2025-05-08]{\normalsize May 08th, 2025}
\maketitle


\part{Faithfulness}
\frame{\partpage}

\begin{frame}{Faithfulness}
A simple iSCM is called $\sigma$-faithful ($d$-faithful) if each conditional independence in its Markov kernel
is due to a $\sigma$-separation ($d$-separation).
\begin{Def}[\ref{def:faithfulness}]
  Let $M = \lt \Sin, \Send, \Srnd, \CX, P, f \rt$ be a simple iSCM with causal graph $G(M)$
  and Markov kernel $P_{M}(X_V \given \doit(X_J))$.
  $M$ is called \emph{$\sigma$-faithful} if 
  for all $A, B, C \ins J \cup V$:
  \begin{equation}\label{slides:eq:sigma_faithful}
    A \Perp^\sigma_{G(M)} B \given C \impliedby X_A \Indep_{P_{M}(X_V \given \doit(\XJ))} X_B \given X_C.
  \end{equation}
  $M$ is called \emph{$d$-faithful} if
  for all $A, B, C \ins J \cup V$:
  \begin{equation*}
    A \Perp^d_{G(M)} B \given C \impliedby X_A \Indep_{P_{M}(X_V \given \doit(\XJ))} X_B \given X_C.
  \end{equation*}
%  If one wants to make the implicit dependence on $J$ more explicit one can equivalently also write:
%  $$A \Perp^\sigma_{G(M)} J \cup B \given C \impliedby X_A \Indep_{P_{M}(X_V,X_W \given \doit(\XJ))} X_J,X_B \given X_C.$$
  For a subset $O \subseteq V$, we say that $M$ is \emph{$\sigma$-faithful w.r.t.\ $O$} if
  \eref{slides:eq:sigma_faithful} holds for all $A, B, C \ins J \cup O$, and define
  \emph{$d$-faithful w.r.t.\ $O$} analogously.
\end{Def}
\end{frame}

\begin{frame}{Faithfulness Violations}
These definitions behave properly under marginalization:
\begin{Rem}
A simple iSCM $M$ is $\sigma$/$d$-faithful w.r.t.\ $O$ if and only if
its marginalization $M_O$ is $\sigma$/$d$-faithful.
\end{Rem}

\bigskip
Faithfulness may fail for various reasons:
\begin{itemize}
  \item Deterministic relationships may lead to additional conditional independences, but are not exploited by the Markov property;
  \item Effects may cancel out;
  \item If cycles are present, and (i) all variables are discrete, or (ii) interactions are linear, $\sigma$-faithfulness may fail;
  \item If cycles are present, and the system is ``perfectly adapting''.
  \item \dots
\end{itemize}
\end{frame}

\begin{frame}{Faithfulness Violation Examples}
\begin{Eg}[Determinism]
  Take an SCM with three endogenous variables $X,Y,Z$ and two exogenous random variables $U,W$, with structural equations
  $$X = 5, \qquad Y = X + U, \qquad Z = X + W.$$
  Then $Y \nPerp Z$ but $Y \Indep Z$.
  This simple (even acyclic) SCM is not faithful due to $X$ being constant.
\end{Eg}

\begin{Eg}[Canceling Effects]
  Take an SCM with three endogenous variables $X,Y,Z$ and three exogenous random variables $W_X,W_Y,W_Z$, with structural equations
  $$X = W_X, \qquad Y = X + W_Y, \qquad Z = Y - X + W_Z.$$
  Then $X \nPerp Z$ but $X \Indep Z$.
\end{Eg}
\end{frame}

\begin{frame}{$\sigma$-Faithfulness Violation}
In certain special cases, the global Markov property in terms of $d$-separation even holds for simple SCMs.
\begin{Prp}[\ref{prp:when_d-separation_holds}]
  Let $M = \lt \Sin, \Send, \Srnd, \CX, P, f \rt$ be a simple SCM with causal graph $G(M)$
  and Markov kernel $P_{M}(X_V \given \doit(X_J))$.
  If $J = \emptyset$ and:
  \begin{enumerate}
    \item $M$ is acyclic, or
    \item all spaces $\CX_v$ with $v \in V$ are discrete, or
%    \item the causal mechanism $f$ is affine and the exogenous distribution has a density w.r.t.\ Lebesgue measure, or
    \item %the causal mechanism $f$ is affine and the exogenous distribution has a density w.r.t.\ Lebesgue measure 
      $M$ is essentially linear over $\RN$ with $\CX_v = \mathbb{R}^{d_v}$ for $d_v \in \mathbb{N}$ for $v \in V$ with
      causal mechanisms of the form
      %$$f(x_V,x_W) = B x_V + c(x_W) \qquad M\as$$
      %where $B : \CX_V \to \CX_V$ is linear and $c : \CX_W \to \CX_V$ is measurable.
      %Since $\CX_V = \bigoplus_{v \in V} \CX_v$, we can also write this as
      $ f_v(x_W,x_V) = \sum_{u\in V} B_{vu} x_u + c_v(x_W) \ M\as$
      where $B : \CX_V \to \CX_V$ is linear and $c : \CX_W \to \CX_V$ is measurable,
      %(that is, $B(x_W)$ is constant).
%      for $v \in V$, where $B_{vu} : \CX_u \to \CX_v$ is a linear mapping for $u,v \in V$ and $c_v : \CX_W \to \CX_v$ is a measurable function for $v \in V$,
  \end{enumerate} 
  then for all $A, B, C \ins J \cup V$:
  $$A \Perp^d_{G(M)} B \given C \implies X_A \Indep_{P_{M}(X_V \given \doit(\XJ))} X_B \given X_C.$$
\end{Prp}
\end{frame}

\begin{frame}{Example: $\sigma$-faithfulness violation}
  \begin{Eg}\small
  \begin{minipage}{0.4\textwidth}\centering
  \scalebox{0.9}{\begin{tikzpicture}
      \node[ndout] (X1) at (0,0) {$X_1$};
      \node[ndout] (X2) at (1.5,0) {$X_2$};
      \node[ndout] (X3) at (1.5,1.5) {$X_3$};
      \node[ndout] (X4) at (0,1.5) {$X_4$};
      \node[ndlat] (W1) at (-1,-1) {$W_1$};
      \node[ndlat] (W2) at (2.5,-1) {$W_2$};
      \node[ndlat] (W3) at (2.5,2.5) {$W_3$};
      \node[ndlat] (W4) at (-1,2.5) {$W_4$};
      \draw[arout] (X1) edge (X2);
      \draw[arout] (X2) edge (X3);
      \draw[arout] (X3) edge (X4);
      \draw[arout] (X4) edge (X1);
      \draw[arout] (W1) edge (X1);
      \draw[arout] (W2) edge (X2);
      \draw[arout] (W3) edge (X3);
      \draw[arout] (W4) edge (X4);
  \end{tikzpicture}}
    \begin{align*}
      X_1 &\idPerp_{G^+(M)} X_3 \given X_2, X_4\\
      X_1 &\nisPerp_{G^+(M)} X_3 \given X_2, X_4\\
      X_1 &\Indep_{P_M} X_3 \given X_2, X_4
    \end{align*}
  \end{minipage}\quad
  \begin{minipage}{0.55\textwidth}
SCM $M$ with structural equations:
  \begin{align*}
    X_1&=X_4 + W_1 \\
    X_2&=X_1 + W_2 \\
    X_3&=X_2 + W_3 \\
    X_4&=\tfrac{1}{2} (X_3 + W_4).
  \end{align*}
We can solve for $X$ in terms of $W$:
  \begin{align*}
    X_1&=2 W_1 + W_2 + W_3 + W_4 \\
    X_2&=2 W_1 + 2 W_2 + W_3 + W_4 \\
    X_3&=2 W_1 + 2 W_2 + 2 W_3 + W_4 \\
    X_4&=W_1 + W_2 + W_3 + W_4
  \end{align*}
This SCM is not $\sigma$-faithful.
  \end{minipage}
\end{Eg}
\end{frame}

\part{Causal Relations}
\frame{\partpage}

\begin{frame}{Remember $\doit(I_B)$ (graphical)}
\begin{Def}[\ref{def:cdmg_intervention_nodes}]
  Let $G=(J,V,E,L)$ be a CDMG and $B \ins J \cup V$ a subset of nodes.
  The \emph{extended CDMG} of $G$ w.r.t.\ nodes $B \ins J \cup V$  and corresponding \emph{intervention nodes} 
  $I_B = \{I_b \,|\, b \in B\}$ {\color{red}with $I_j:=j$ for $j \in B \cap J$},
  is the CDMG: 
  \[ G_{\doit(I_B)}:=(J_{\doit(I_B)},V_{\doit(I_B)},E_{\doit(I_B)},L_{\doit(I_B)}),\]
  %\[G(V|\doit(J \dcup I_W)):=G_{\doit(I_W)}=(J_{\doit(I_W)},V_{\doit(I_W)},E_{\doit(I_W)},L_{\doit(I_W)}),\]
  where:
  \begin{enumerate}
      \item $J_{\doit(I_B)}:= J\, \dot{\cup} \, \lC I_b\,|\, b \in B {\color{red}\sm J} \rC$, %not B \sm J, makes weak and hard interventions commute
      \item $V_{\doit(I_B)}:= V$,
      \item $E_{\doit(I_B)}:= E\, \dot{\cup} \, \lC I_b \tuh b \,|\, b \in B\sm J \rC$,
      \item $L_{\doit(I_B)}:= L$.
  \end{enumerate}
\end{Def}
\end{frame}

\begin{frame}{Remember $\doit(I_B)$ (iSCM)}
  \begin{Def}[\ref{def:scm_intervention_variables}]
  Let $M = \lt \Sin, \Send, \Srnd, \CX, P, f \rt$ be an iSCM, and $B \subseteq V \cup J$.
    For each $b \in B \cap V$, we introduce an additional \emph{intervention variable} $I_b$ as exogenous input variable; \alert{for $j \in B \cap J$, we just set $I_j := j$}.
  We denote $I_B := (I_b)_{b \in B}$.
  We define an extended iSCM 
  \vskip-2mm
  $$\textstyle M_{\doit(I_B)} := \lt \Sin \cup I_B, \Send, \Srnd, \CX \times \prod_{b \in B \cap V} \CX_{I_b}, P, \tilde{f} \rt$$
  with $\CX_{I_b} := \CX_b \,\dot\cup\, \{\star\}$, with causal mechanism $\tilde{f}$ with components
  \vskip-6mm
  \begin{align*}
    v \in V \cap B: \qquad &\tilde{f}_v(x_{V \cup J \cup W \cup I_B}) := \begin{cases}
      f_v(x_{V \cup J \cup W}) & x_{I_v} = \star \\
      x_{I_v} & x_{I_v} \in \CX_v
    \end{cases}\\
    v \in V \sm B: \qquad & \tilde{f}_v(x_{V \cup J \cup W \cup I_B}) := f_v(x_{V \cup J \cup W})
  \end{align*}
  \vskip-2mm
  For $b \in B \cap V$, $x_{I_b} = \star$ encodes that there is no intervention on $b$, while $x_{I_b} \ne \star$ encodes that the hard intervention $\doit(X_b = x_{I_b})$ is performed.
\end{Def}
Purpose: jointly encode $\{M_{\doit(X_C)} : C \subseteq B\}$ into a single iSCM $M_{\doit(I_B)}$.
\end{frame}

\begin{frame}{Outline}
  In this chapter, we formalize the notions:
  \begin{enumerate}
    \item $a$ causes $b$
    \item $a$ directly causes $b$
    \item $a$ and $b$ have a common cause
    \item $a$ and $b$ are confounded
  \end{enumerate}

  \bigskip
  We distinguish different versions of these notions:
  \begin{enumerate}
    \item Graphical
    \item Counterfactual (``rung-3'') / Potential outcome
    \item Interventional (``rung-2'')
  \end{enumerate}

  \bigskip
  Analogous graphical and interventional notions apply to CBNs.
\end{frame}


\begin{frame}{Causal relations (graphical notion)}
\begin{Def}[\ref{def:graph_cause_of}]
  Let $G$ be a CDMG with input nodes $J$ and output nodes $V$.
  Let $a \in J \cup V$ and $b \in V$.
  If $a \notin \Anc^G(b)$ then we say that \emph{$a$ does not cause $b$ according to $G$}; otherwise, we say that \emph{$a$ causes $b$ according to $G$}.
\end{Def}

\bigskip
  Remarks:
  \begin{itemize}
    \item Marginalization of a graph preserves its causal relations (see also Remark~\ref{rem:marg_preserves_ancestors_bifurcations_acyclicity}).
    \item $a$ causes $b$ according to $G$ \\
          $\iff$ $a \in \Anc^G(b)$ \\
          $\iff$ $a \tuh \dots \tuh b$ in $G$ \\
          $\iff$ $a \tuh b$ is present in in $G^{\sm(V\sm\{a,b\})}$ \\
          $\iff$ $a$ causes $b$ according to $G_{\doit(a)}$ \\
          $\iff$ $I_a$ causes $b$ according to $G_{\doit(I_a)}$.
  \end{itemize}
\end{frame}

\begin{frame}{Causal relations (graphical notion)}
\begin{Prp}[\ref{prp:graph_cause_of}]
  Let $G$ be a CDMG with input nodes $J$ and output nodes $V$.
  Let $b \in V$.
  For $a \in J \cup V$:
    \begin{equation*}
      a \not\in \Anc^G(b) \iff b \Perp^\sigma_{G_{\doit(I_a)}} I_a \mid J \sm \{I_a\}.
    \end{equation*}
\end{Prp}

\begin{Rem}
  We can explicitly distinguish the cases $a \in J$ and $a \in V$ as follows:
  $$b \Perp^\sigma_{G_{\doit(I_a)}} I_a \mid J \sm \{I_a\} \iff \begin{cases}
    b \Perp^\sigma_{G} a \mid J \setminus \{a\} & a \in J \\
    b \Perp^\sigma_{G_{\doit(I_a)}} I_a \mid J & a \in V.
  \end{cases}$$
\end{Rem}
\end{frame}


\begin{frame}[t]{Causal relations (interventional notion)}
\begin{Def}[\ref{def:scm_rung2_cause_of}]
  Let $M = \lt \Sin, \Send, \Srnd, \CX, P, f \rt$ be a simple iSCM.
  Let $a \in J \cup V$ and $b \in V$.
  If
    \begin{equation*}
      X_b \Indep_{P_{M_{\doit(I_a)}}} X_{I_a} \given X_{J \sm \{I_a\}},
    \end{equation*}
  we say that \emph{$a$ is not a rung-2 cause of $b$ according to $M$}.
  Otherwise, we say that \emph{$a$ is a rung-2 cause of $b$ according to $M$}.
\end{Def}

\bigskip
  \only<1>{We can rewrite \eref{eq:scm_rung2_cause_of}, distinguishing the cases $a \in J$ and $a \in V$, as follows:
  $$X_b \Indep_{P_{M_{\doit(I_a)}}} X_{I_a} \given X_{J \sm \{I_a\}} \iff \begin{cases}
    X_b \Indep_{P_{M}} X_a \given X_{J \sm \{a\}} & a \in J, \\
    X_b \Indep_{P_{M_{\doit(I_a)}}} X_{I_a} \given X_J & a \in V.
  \end{cases}$$}
  \only<2>{\small Alternatively, we can write \eref{eq:scm_rung2_cause_of} as:
    $$P_{M}(X_b \given \doit(X_{J \cup \{a\}})) = P_{M}(X_b \given \doit(X_{J\sm\{a\}})),$$
  or, distinguishing the two cases, as:
    $$\begin{cases}
      P_{M}(X_b \given \doit(X_J)) = P_{M}(X_b \given \doit(X_{J\setminus \{a\}}),\cancel{\doit(X_a)}) & a \in J, \\
      P_{M}(X_b \given \doit(X_J), \doit(X_a)) = P_{M}(X_b \given \doit(X_J)) & a \in V.
    \end{cases}$$
    }
\end{frame}

\begin{frame}{Causal relations (interventional notion)}
  Remarks:
  \begin{itemize}
    \item Marginalization of an iSCM preserves the rung-2 causal relations between the remaining variables.
    \item Rung-2 causal relations are invariant under interventional equivalence of iSCMs.
    \item We treat the two cases $a \in V$ and $a \in J$ differently, and one may wonder why we chose to work with $M_{\doit(I_a)}$ rather than $M_{\doit(a)}$. 
      The reason is: the observable Markov kernel may provide additional information beyond that provided by the interventional Markov kernels (see Example~\ref{eg:scm_cause_of_include_observational}).
  \end{itemize}
\end{frame}

\begin{frame}{Causal relations (interventional notion)}
  \begin{Eg}[\ref{eg:scm_cause_of_include_observational}]
  Consider the SCM $M$ with endogenous variables $X_a, X_b \in \{-1,+1\}$, exogenous random variable $X_w \in \{-1,1\}$, and structural equations:
  \begin{align*}
    X_a &= X_w, \\
    X_b &= X_a X_w,
  \end{align*}
  where $X_w \sim \C{U}(\{-1,1\})$.

  Note that 
    $$P_M(X_b = 1 \given \doit(X_a = x_a)) = \tfrac{1}{2} \quad\text{for $x_a \in \{-1,1\}$},$$
  while
    $$P_M(X_b = 1) = 1.$$
  Hence, $a$ is a rung-2 cause of $b$ according to $M$\\
  (but not according to $M_{\doit(a)}$!).
\end{Eg}
\end{frame}

\begin{frame}{Causal relations (counterfactual notion)}
\begin{Def}[\ref{def:scm_rung3_cause_of}]
  Let $M = \lt \Sin, \Send, \Srnd, \CX, P, f \rt$ be a simple iSCM.
  Let $a \in J \cup V$ and $b \in V$.
  We say that \emph{$a$ is not a rung-3 cause of $b$ according to $M$} if:
    \begin{equation*}\begin{split}
      & \forall x_{J\sm\{I_a\}} \in \CX_{J\sm\{I_a\}}, \forall x_{I_a},x_{I_{a'}} \in \CX_a: \\
      & P_{(M_{\doit(X_{J\sm\{I_a\}}=x_{J\sm\{I_a\}}),\doit(I_a)})^\twin}(X_b = X_{b'} \given \doit(X_{I_a}=x_{I_a}),\doit(X_{I_{a'}}=x_{I_{a'}})) = 1.
    \end{split}\end{equation*}
  Otherwise, we say that \emph{$a$ is a rung-3 cause of $b$ according to $M$}.
\end{Def}
  We can rewrite the condition for $a \in V$ as:
%    $$P_{M}(X_b \given \doit(X_J), \doit(X_a)) \ne P_{M}(X_b \given \doit(X_J))$$
    $$\forall x_J \in \CX_J, \forall x_a \in \CX_a: \quad P_{((M_{\doit(X_J=x_J)})^\twin)_{\doit(a)}}(X_b = X_{b'} \given \doit(X_a=x_a)) \ne 1;$$
  and for $a \in J$ as:
%    $$P_{M}(X_b \given \doit(X_J)) \ne P_{M}(X_b \given \doit(X_{J\setminus \{a\}}),\cancel{\doit(X_a)}).$$
    \begin{equation*}\begin{split}
      &\forall x_{J\sm\{a\}} \in \CX_{J\sm\{a\}}, \forall x_a,x_a' \in \CX_a: \\
      &P_{(M_{\doit(X_{J\sm\{a\}}=x_{J\sm\{a\}})})^\twin}(X_b = X_{b'} \given \doit(X_a=x_a),\doit(X_{a'}=x_a')) \ne 1.
    \end{split}\end{equation*}
\end{frame}

\begin{frame}{Causal relations (counterfactual notion)}
  \begin{itemize}
    \item Marginalization of an iSCM preserves the rung-3 causal relations of the remaining variables.
    \item Rung-3 causal relations are invariant under counterfactual equivalence of iSCMs.
    \item The presence of a rung-3 causal relation does not necessarily imply the presence of a rung-2 causal relation (Example~\ref{eg:scm_cause_of_rung2_vs_rung3}).
    Intuitively, the causal effect can go unnoticed if it averages out.
  \item The presence of a causal relation according to the causal graph does not necessarily imply the presence of a rung-3 causal relation (Example~\ref{eg:scm_cause_of_graphical_vs_rung3}).
Intuitively, different causal paths can cancel out.
  \end{itemize}
\end{frame}

\begin{frame}{Causal relations: rung-2 vs.\ rung-3}
  \begin{Eg}[\ref{eg:scm_cause_of_rung2_vs_rung3}]
  Consider the iSCM $M$ with exogenous input variable $X_a \in \{-1,+1\}$, endogenous variable $X_b \in \{-1,+1\}$, exogenous random variable $X_w \in \{-1,1\}$, and structural equation:
  $$X_b = X_a X_w,$$
  where $X_w \sim \C{U}(\{-1,1\})$.
  Note that $a$ is a rung-3 cause of $b$ according to $M$, but $a$ is not a rung-2 cause of $b$ according to $M$.
  Indeed, due to the uniform distribution of $X_w$ the dependence of $b$ on $a$ averages out.
\end{Eg}
  For example, a drug can have a beneficial effect on all males and an adversary effect on all females, and still have \emph{no effect} on the entire population (if the positive effect exactly cancels out the negative effect on average).
\end{frame}

\begin{frame}{Causal relations: graphical vs.\ rung-3}
  \begin{Eg}[\ref{eg:scm_cause_of_graphical_vs_rung3}]
  Consider the iSCM $M$ with exogenous input variable $X_a \in \RN$, endogenous variables $X_c, X_d, X_b \in \RN$, exogenous random variable $X_w \in \RN$, and structural equations
  \begin{align*}
    X_c &= X_a, \\
    X_d &= -X_a \\
    X_b &= X_c + X_d + X_w 
  \end{align*}
  with $X_w \sim \C{N}(0,1)$.
  Then $a$ causes $b$ according to $G(M)$, but $a$ is not a rung-3 cause of $b$ according to $M$.
\end{Eg}
\end{frame}

\begin{frame}{Causal relations (potential outcome notion)}
We can also use potential outcomes to formulate the notion of causation.
\begin{Def}[\ref{def:scm_po_cause_of}]
  Let $M = \lt \Sin, \Send, \Srnd, \CX, P, f \rt$ be a simple iSCM.
  Let $a \in J \cup V$ and $b \in V$.
  Let $g$ be a solution function of $M_{\doit(I_a)}$ and let $X_W \sim P(X_W)$ be a random variable.
  We say that \emph{$a$ is not a potential-outcome cause of $b$ according to $M$} if:
    \begin{equation*}\begin{split}
   & \forall x_{J\sm\{I_a\}} \in \CX_{J\sm\{I_a\}}, \forall x_{I_a},x_{I_a}' \in \CX_{I_a}: \\
      & g_b(x_{I_a},x_{J\sm\{I_a\}},X_W) = g_b(x_{I_a}',x_{J\sm\{I_a\}},X_W) \text{ a.s..}
    \end{split}\end{equation*}
  Otherwise, \emph{$a$ is a potential-outcome cause of $b$ according to $M$}.
\end{Def}
\small
Using the notation of potential outcomes, we can also write:
  $$\forall x_{J\sm\{I_a\}} \in \CX_{J\sm\{I_a\}}, \forall x_{I_a},x_{I_a}' \in \CX_{I_a}: \qquad X^{\doit(x_{I_a},x_{J\sm\{I_a\}})}_b = X^{\doit(x_{I_a}',x_{J\sm\{I_a\}})}_b \text{ a.s.}$$
  \alert{It is very important to keep in mind that we must use the same $X_W$ to induce $X^{\doit(x_{I_a},x_{J\sm\{I_a\}})}_b := g_b(x_{I_a},x_{J\sm\{I_a\}},X_W)$ and $X^{\doit(x_{I_a}',x_{J\sm\{I_a\}})}_b := g_b(x_{I_a}',x_{J\sm\{I_a\}},X_W)$.}
\end{frame}

\begin{frame}{Causal relations: counterfactual vs.\ potential outcome}
At second sight, it turns out that the potential-outcome notion and the counterfactual notion coincide.
  \begin{Prp}[\ref{prp:scm_cause_of_rung3_po}]
  Let $M = \lt \Sin, \Send, \Srnd, \CX, P, f \rt$ be a simple iSCM.
  Let $a \in J \cup V$ and $b \in V$.
  Then $a$ is a rung-3 cause of $b$ if and only if $a$ is a potential-outcome cause of $b$.
\end{Prp}
The potential-outcome notion has the advantage that it is considerably simpler to formulate: it leads to shorter expressions than the ones involving the twinning operation.
\end{frame}

\begin{frame}{Hierarchy of causal relations}

\begin{Thm}[\ref{thm:scm_causes}]
  Let $M = \lt J, V, W, \CX, P, f \rt$ be a simple iSCM with causal graph $G(M)$.
  Let $a \in J \cup V$, $b \in V$. Consider the following statements:
  \begin{enumerate}
    \item[(i)] $a$ does not cause $b$ according to $G(M)$;
    \item[(ii)] $a$ is not a rung-3 cause of $b$ according to $M$;
    \item[(iii)] $a$ is not a rung-2 cause of $b$ according to $M$.
  \end{enumerate}
  Then (i) $\implies$ (ii) $\implies$ (iii). If $M_{\doit(I_a)}$ is faithful, also (iii) $\implies$ (i).
\end{Thm}
\end{frame}

\begin{frame}{Direct causal relations}
  \begin{itemize}
    \item Essentially, ``$a$ causes $b$ directly'' means that $a$ causes $b$ even when performing a hard intervenion on all other endogenous variables.
    \item Similarly to ``$a$ causes $b$'', we also distinguish different versions of ``$a$ causes $b$ directly''.
    \item One should keep in mind that the notion of direct causation is always relative to some set of variables.
    \item In particular, this property is not necessarily preserved under marginalization!
  \end{itemize}
\end{frame}

\begin{frame}{Direct causal relations (graphical notion)}
\begin{Def}[\ref{def:graph_submodel}]
  Let $G = \lt J, V, E, L \rt$ be a CDMG with input nodes $J$ and output nodes $V$.
  For $A \subseteq V$, we define the \emph{submodel} of $G$ on $A$ as $G^{[A]} := G^{\doit(V\sm A)} = \lt J \cup (V \sm A), A, \lC v \tuh a \in E \given v \in V, a \in A \rC, \lC a \huh a' \in L \given a,a' \in A \rC \rt$.
\end{Def}
\begin{Def}[\ref{def:graph_direct_cause_of}]
  Let $G$ be a CDMG with input nodes $J$ and output nodes $V$.
  Let $a \in J \cup V$ and $b \in V$.
  We say that \emph{$a$ is a direct cause of $b$ (w.r.t.\ $V \cup J$) according to $G$} if $a$ causes $b$ according to $G^{[\{a,b\}\cap V]}$.
\end{Def}
  This holds if and only if $a \in \Pa^{G}(b)$
  (so this provides the causal interpretation of directed edges in a causal graph).
\end{frame}

\begin{frame}{Direct causal relations (interventional \& counterfactual notion)}
\begin{Def}[\ref{def:scm_rung2_direct_cause_of}]
  Let $M = \lt \Sin, \Send, \Srnd, \CX, P, f \rt$ be a simple iSCM.
  Let $a \in J \cup V$ and $b \in V$. 
  We say that \emph{$a$ is a rung-2 direct cause of $b$ (w.r.t.\ $J \cup V$) according to $M$} if and only if $a$ is a rung-2 cause of $b$ according to $M^{[\{a,b\}\cap V]}$.
\end{Def}
  Rung-2 direct causal relations are invariant under interventional equivalence of iSCMs.
\begin{Def}[\ref{def:scm_rung3_direct_cause_of}]
  Let $M = \lt \Sin, \Send, \Srnd, \CX, P, f \rt$ be a simple iSCM.
  Let $a \in J \cup V$ and $b \in V$. 
  We say that \emph{$a$ is a rung-3 direct cause of $b$ (w.r.t.\ $J \cup V$) according to $M$} if and only if $a$ is a rung-3 cause of $b$ according to $M^{[\{a,b\}\cap V]}$.
  %  $= M_{\doit(V\sm \{a,b\})}$
\end{Def}
  Rung-3 direct causal relations are invariant under counterfactual equivalence of iSCMs.
\end{frame}

\begin{frame}{Direct causal relations (hierarchy)}
\begin{Cor}[\ref{cor:scm_direct_cause_of}]
  Let $M = \lt \Sin, \Send, \Srnd, \CX, P, f \rt$ be a simple iSCM with causal graph $G(M)$.
  Let $a \in V \cup J$ and $b \in V$.
  Consider the following statements:
  \begin{enumerate}
    \item[(i)] $a$ is not a direct cause of $b$ (w.r.t.\ $V \cup J$) according to $G(M)$;
    \item[(ii)] $a$ is not a rung-3 direct cause of $b$ (w.r.t.\ $V \cup J$) according to $M$;
    \item[(iii)] $a$ is not a rung-2 direct cause of $b$ (w.r.t.\ $V \cup J$) according to $M$.
  \end{enumerate}
  (i) $\implies$ (ii) $\implies$ (iii). If $M^{[\{a,b\}\cap V]}_{\doit(I_a)}$ is faithful, also (iii) $\implies$ (i).
\end{Cor}
\end{frame}

\begin{frame}{Common causes}
  \begin{itemize}
    \item Essentially, ``$c$ is a common cause of $a,b$'' means that $c$ causes $b$ even when performing a hard intervention on $a$, and $c$ causes $a$ even when performing a hard intervention on $b$.
    \item Equivalently, it means that $c$ is a direct cause of both $a$ and $b$ with respect to $\{a,b,c\}$.
    \item We distinguish different versions of ``$a$ and $b$ have common cause $c$''.
  \end{itemize}
\end{frame}

\begin{frame}{Common causes (various notions)}
\begin{Def}[\ref{def:graph_common_cause_of}]
  Let $G$ be a CDMG with input nodes $J$ and output nodes $V$.
  Let $a, b \in V$ and $c \in V \cup J$ such that $a,b,c$ are distinct.
  We say that \emph{$c$ is a common cause of $a$ and $b$ according to $G$} if $c$ causes $a$ according to $G_{\doit(b)}$ and $c$ causes $b$ according to $G_{\doit(a)}$.
\end{Def}
\begin{Def}[\ref{def:scm_rung2_common_cause_of}, \ref{def:scm_rung3_common_cause_of}]
  Let $M = \lt \Sin, \Send, \Srnd, \CX, P, f \rt$ be a simple iSCM.
  Let $a, b \in V$ and $c \in V \cup J$ such that $a,b,c$ are distinct.
  We say that:
  \begin{itemize}
    \item \emph{$c$ is a rung-2 common cause of $a,b$ according to $M$} if and only if $c$ is a rung-2 cause of $a$ according to $M_{\doit(b)}$ and $c$ is a rung-2 cause of $b$ according to $M_{\doit(a)}$;
    \item \emph{$c$ is a rung-3 common cause of $a,b$ according to $M$} if and only if $c$ is a rung-3 cause of $a$ according to $M_{\doit(b)}$ and $c$ is a rung-3 cause of $b$ according to $M_{\doit(a)}$.
  \end{itemize}
\end{Def}
\end{frame}

\begin{frame}{Common causes (properties)}
  \begin{itemize}
    \item
     $a$ and $b$ have a common cause according to $G$\\
     $\iff$ $a \hut c$ and $c \tuh b$ are both present in $G^{\sm(V \sm \{a,b,c\})}$ \\
     $\iff$ there exists a bifurcation with source $c$ in $G$ between $a$ and $b$.
   \item
     The rung-2 (resp.\ rung-3) notion of having a common cause is invariant under interventional (resp.\ counterfactual) equivalence of iSCMs.
  \end{itemize}
\begin{Def}[Part of Definition~\ref{def:walks}]
    Let $G = \lt J, V, E, L \rt$ be a CDMG. 
    A \emph{bifurcation} between $v$ and $w$ in $G$ is a walk of the form:
        \[v=v_0 \hut v_1 \hut  \cdots \hut v_{k-1} \hus v_k  \tuh \cdots \tuh v_{n-1} \tuh v_n=w,\]
   %             \[v=v_0 \hut v_1 \hut  \cdots \hut v_{k-1} \huh v_k  \tuh \cdots \tuh v_{n-1} \tuh v_n=w,\]
        such that $v \ne w$, the walk contains both endnodes exactly once, 
        every node has at most one arrowhead pointing towards it, and 
        both endnodes have exactly one arrowhead pointing towards them. 
          If the edge $v_{k-1} \hus v_k$ is directed ($v_{k-1} \hut v_k$) then we say that the bifurcation has \emph{source} $v_{k}$.
  \end{Def}
\end{frame}

\begin{frame}{Common causes (interventional \& counterfactual notions)}

\begin{Cor}[\ref{cor:scm_common_cause_of}]
  Let $M = \lt \Sin, \Send, \Srnd, \CX, P, f \rt$ be a simple iSCM with causal graph $G(M)$.
  Let $a, b \in V$ and $c \in V \cup J$ such that $a,b,c$ are distinct.
  Consider the following statements:
  \begin{enumerate}
    \item[(i)] $c$ is not a common cause of $a,b$ according to $G(M)$;
    \item[(ii)] $c$ is not a rung-3 common cause of $a,b$ according to $M$;
    \item[(iii)] $c$ is not a rung-2 common cause of $a,b$ according to $M$.
  \end{enumerate}
  (i) $\implies$ (ii) $\implies$ (iii). If $M_{\doit(b),\doit(I_c)}$ and $M_{\doit(a),\doit(I_c)}$ are both faithful, also (iii) $\implies$ (i).
\end{Cor}
\end{frame}

\begin{frame}{Confounded}
  \begin{itemize}
    \item The notion of ``confounded'' (or ``spurious dependence'') is trickiest to define (especially in case of cycles).
    \item Essentially, we say that two endogenous variables are confounded if their joint distribution is not entirely explained by their mutual causal effects.
  \end{itemize}
\end{frame}

\begin{frame}{Confounding (graphical)}
\begin{Def}[\ref{def:graph_unconfounded}]
  Let $G$ be a CDMG with input nodes $J$ and output nodes $V$.
  Let $a,b \in V$ be distinct output nodes.
  If there exists a bifurcation between $a$ and $b$ in $G$, either without source or with source $c \in V$, then we say that \emph{$a$ and $b$ are confounded according to $G$}.
  Otherwise, we say that \emph{$a$ and $b$ are unconfounded according to $G$}.
\end{Def}
  \begin{itemize}
    \item The reason we exclude bifurcations with as source an exogenous input node is that these will not lead to ``spurious dependences'' in the ``strata'' of the Markov kernel.
    \item  Marginalization of a graph preserves confoundedness (see also Remark~\ref{rem:marg_preserves_ancestors_bifurcations_acyclicity}).
    \item In particular, $a$ and $b$ are confounded according to $G$ if and only if $a \huh b$ %or $a \hut j \tuh b$ with $j \in J$ 
  is present in $G^{\sm (V \sm \{a,b\})}$.
%  Proposition~\ref{prp:bifurcations_alternative} gives yet another equivalent graphical formulation of this notion.
  \end{itemize}
\end{frame}


\begin{frame}{Confounding (graphical notion)}
\begin{Prp}[\ref{prp:graph_unconfounded2}]
%  Let $M = \lt \Sin, \Send, \Srnd, \CX, P, f \rt$ be a simple iSCM with causal graph $G(M)$.
%  Let $a, b \in V$ with $a \ne b$.
%  If $a$ and $b$ are unconfounded according to $G(M)$ then 
%  $$a \Perp^\sigma_{G(M^\twin_{\doit(a',b)})} b' \given J, J', a', b.$$
  Let $G^+$ be a CDMG with input nodes $J$ and output nodes $V \cup W$ such that the nodes in $W$ have no parents in $G^+$.
  Let $a,b \in V$ be distinct output nodes in $V$.
  If $a$ and $b$ are unconfounded according to $G := (G^+)^{\sm W}$, then
  $$a \isPerp_{(G^+)^{\twin(J \cup V)}_{\doit(a',b)}} b' \given J, J', a', b.$$
  The converse holds if each node in $V$ has at least one parent in $W$ in $G^+$.
\end{Prp}
\end{frame}

\begin{frame}{The intuition behind \dots}
  \begin{tikzpicture}
    \begin{scope}[xshift=-3cm]
      \node[ndout] (X) at (0,0) {$a$};
      \node[ndout] (Y) at (0,-1.5) {$b$};
      \draw[arout] (X) to (Y);
      \draw[arout,bend left] (Y) to (X);
      \draw[arlat,bend left] (X) to (Y);
      \node at (0,-3) {$(G^+)^{\sm W}$};
    \end{scope}
    \begin{scope}
      \node[ndout] (X) at (0,0) {$a$};
      \node[ndout] (Y) at (0,-1.5) {$b$};
      \node[ndlat] (W1) at (1,0.5) {$w_a$};
      \node[ndlat] (W) at (1,-0.75) {$w$};
      \node[ndlat] (W3) at (1,-2) {$w_b$};
      \draw[arout] (X) to (Y);
      \draw[arout,bend left] (Y) to (X);
      \draw[arout] (W1) to (X);
      \draw[arout] (W) to (X);
      \draw[arout] (W) to (Y);
      \draw[arout] (W3) to (Y);
      \node at (0.5,-3) {$G^+$};
    \end{scope}
    \begin{scope}[xshift=5cm]
      \node[ndout] (X) at (0,0) {$a$};
      \node[ndint] (Y) at (0,-1.5) {$b$};
      \node[ndlat] (W1) at (1,0.5) {$w_a$};
      \node[ndlat] (W) at (1,-0.75) {$w$};
      \node[ndlat] (W3) at (1,-2) {$w_b$};
%      \draw[arout] (X) to (Y);
      \draw[arout,bend left] (Y) to (X);
      \draw[arout] (W1) to (X);
      \draw[arout] (W) to (X);
%      \draw[arout] (W) to (Y);
%      \draw[arout] (W3) to (Y);
      \node[ndint] (X2) at (2,0) {$a'$};
      \node[ndout] (Y2) at (2,-1.5) {$b'$};
      \draw[arout] (X2) to (Y2);
%      \draw[arout,bend right] (Y2) to (X2);
%      \draw[arout] (W1) to (X2);
%      \draw[arout] (W) to (X2);
      \draw[arout] (W) to (Y2);
      \draw[arout] (W3) to (Y2);
      \draw[dotted] (-1,-2.5) rectangle (1,1.0);
      \draw[dotted] (1,-2.5) rectangle (3,1.0);
      \node at (0,1.3) {\tiny factual};
      \node at (2,1.3) {\tiny counterfactual};
      \node at (1,-3) {$(G^+)^{\twin(J \cup V)}_{\doit(a',b)}$};
    \end{scope}
  \end{tikzpicture}
  Proposition~\ref{prp:graph_unconfounded2} claims that if $a$ and $b$ are unconfounded according to $(G^+)^{\sm W}$ then there is no $\sigma$-open path between $a$ and $b'$ given $a',b$ $(G^+)^{\twin(J \cup V)}_{\doit(a',b)}$.
\end{frame}

\begin{frame}{Confounding (counterfactual notion)}
  
\begin{Def}[\ref{def:scm_rung3_unconfounded}]
  Let $M = \lt \Sin, \Send, \Srnd, \CX, P, f \rt$ be a simple iSCM.
  Let $a, b \in V$ with $a \ne b$.
  Let $g^{\doit(a)}$ be a solution function of $M_{\doit(a)}$, and $g^{\doit(b)}$ a solution function of $M_{\doit(b)}$.
  If for all values $x_J \in \CX_J, x_{J'} \in \CX_J$, $x_{a'} \in \CX_a$, $x_b \in \CX_b$: 
  $$g^{\doit(b)}_a(x_b,x_J,X_W) \Indep g^{\doit(a)}_b(x_{a'},x_{J'},X_W)$$
  with $X_W \sim P$, then we say that \emph{$a$ and $b$ are rung-3 unconfounded according to $M$}.
\end{Def}

  \begin{itemize}
    \item This is equivalent to:
      $$X_a \Indep_{P_{M^\twin_{\doit(a',b)}}} X_{b'} \given X_J, X_{J'}, X_{a'}, X_b.$$
    \item Rung-3 unconfoundedness is invariant under counterfactual equivalence of iSCMs.
    \item Rung-3 unconfoundedness is \emph{not} invariant under interventional equivalence of iSCMs (Example~\ref{eg:scm_cause_of_include_observational}).
  \end{itemize}
\end{frame}

\begin{frame}{Confounding (counterfactual, no cycle)}
  In case $a$ and $b$ are not part of a causal cycle, one can use another formulation. This one occurs more commonly in the literature.
  \begin{Prp}[\ref{prp:scm_rung3_unconfounded_acyclic}]
  Let $M = \lt \Sin, \Send, \Srnd, \CX, P, f \rt$ be a simple iSCM.
  Let $a, b \in V$ with $a \ne b$.
  Suppose that $b$ is not a rung-3 cause of $a$ according to $M$.
  Then $a$ and $b$ are rung-3 unconfounded according to $M$ if and only if
%  for all values $x_J \in \CX_J, x_{J'} \in \CX_J$, $x_{a'} \in \CX_a$: 
%  $$g_a(x_J,X_W) \Indep g^{\doit(a)}_b(x_{a'},x_{J'},X_W),$$
  \begin{equation*}
    \forall x_J,x_{J'} \in \CX_J \ \forall x_{a'} \in \CX_a: g_a(x_J,X_W) \Indep g^{\doit(a)}_b(x_{a'},x_{J'},X_W)
  \end{equation*}
  for $X_W \sim P(X_W)$ (where $g^{\doit(a)}$ denotes a solution function of $M_{\doit(a)}$, and $g$ a solution function of $M$).
\end{Prp}
  \begin{itemize}
    \item The condition is equivalent to:
      $$X_a \Indep_{P_{M^\twin_{\doit(a')}}} X_{b'} \given X_J, X_{J'}, X_{a'}.$$
  \end{itemize}
\end{frame}

\begin{frame}{Exchangeability, ignorability}
  Suppose $J = \emptyset$. 
  Consider the potential outcomes $X_a$ for $M$, $X_a^{\doit(x_b)}$ for $M_{\doit(b)}$ and $X_{b}^{\doit(x_{a})}$ for $M_{\doit(a)}$:
\begin{align*}
  X_a &:= g_a(X_W), \\
  X_{a}^{\doit(x_b)} &:= g^{\doit(b)}_a(x_b,X_W), \qquad x_b \in \CX_b\\
  X_{b}^{\doit(x_a)} &:= g^{\doit(a)}_b(x_a,X_W), \qquad x_a \in \CX_a.
\end{align*}
The assumption that $b$ is not a rung-3 cause of $a$ according to $M$ becomes:
  $$\forall x_b \in \CX_b: \qquad X_a^{\doit(x_b)} = X_a\ \mathrm{a.s.}.$$

The (acyclic) condition for unconfoundedness of $a$ and $b$ according to $M$ becomes:
  $$\forall x_a \in \CX_a : \qquad X_a \Indep X_b^{\doit(x_a)}.$$
This condition is also referred to as \emph{exchangeability} or \emph{ignorability} in the potential outcome literature. 
\end{frame}

\begin{comment}
\begin{Eg}[\ref{eg:scm_cause_of_iv_notimplies_iii}]
  Consider the iSCM $M$ with exogenous input variable $X_a \in \RN$, endogenous variable $X_b \in \RN$, exogenous random variable $X_w \in \RN$, and structural equation
  $$X_b = \ind_{X_a}(X_w)$$
  and exogenous distribution $X_w \sim \C{N}(0,1)$.
  Then for all $x_a \in \RN$, $g(x_a,X_w) = 0$ a.s., so property (iv) in Proposition~\ref{prp:scm_causes} holds.
  However, it is not the case that for almost all $X_w$, we have that $g_b(x_a,X_w) = g_b(x_a',X_w)$ for all $x_a,x_a'$, which means that property (iii) in Proposition~\ref{prp:scm_causes} does not hold.
\end{Eg}

\begin{Eg}[\ref{eg:scm_cause_of_iii_notimplies_ii}]
  Consider the iSCM $M$ with exogenous input variable $X_a \in \{0,1\}$, endogenous variable $X_b \in \{0,1\}$, exogenous random variable $X_w \in \{0,1\}$, and structural equations
  $$X_b = X_a \wedge X_w$$
  with $P(X_w=0) = 1$.
  Then $X_b = 0$ with probability 1, hence property (iii) in Proposition~\ref{prp:scm_causes} holds.
  However, property (ii) in Proposition~\ref{prp:scm_causes} does not hold.
\end{Eg}

  \begin{Eg}[\ref{eg:scm_cause_of_ii_notimplies_i}]
  Consider the iSCM $M$ with exogenous input variable $X_a \in \RN$, endogenous variables $X_c, X_d, X_b \in \RN$, exogenous random variable $X_w \in \RN$, and structural equations
  \begin{align*}
    X_c &= X_a, \\
    X_d &= -X_a \\
    X_b &= X_c + X_d + X_w 
  \end{align*}
  with $X_w \sim \C{N}(0,1)$.
  Then (i) in Proposition~\ref{prp:scm_causes} does not hold, but (ii) in Proposition~\ref{prp:scm_causes} holds.
\end{Eg}
\end{comment}

\begin{frame}{Confounding (interventional notion)}

  Caveat: the following definition applies to the acyclic case, we don't know how to generalize it to the cyclic setting.

  \begin{Def}[\ref{def:no_confounding_bias}]
  Let $M = \lt \Sin, \Send, \Srnd, \CX, P, f \rt$ be a simple iSCM.
  Let $a \ne b \in V$.
%  We say that \emph{$a$ and $b$ have no confounding bias according to $M$} if:
%  $b$ is not a rung-2 cause of $a$ according to $M$ implies
%  \begin{equation}\label{eq:no_confounding_bias}
%    P_{M}(X_b \given \doit(X_a), \doit(X_J)) = P_{M}(X_b \given X_a, \doit(X_J)) \qquad P_M(X_a \given \doit(X_J))\text{-a.s.}.
%  \end{equation}
  Assume $b$ is not a rung-2 cause of $a$ according to $M$.
  Then we say that \emph{$a$ and $b$ have no confounding bias according to $M$} if:
  \begin{equation*}
    P_{M}(X_b \given \doit(X_a), \doit(X_J)) = P_{M}(X_b \given X_a, \doit(X_J)) \qquad P_M(X_a \given \doit(X_J))\text{-a.s.}.
  \end{equation*}
\end{Def}
  \begin{itemize}
    \item We can write the condition equivalently as:
      $$X_b \Indep_{P_{M_{\doit(I_a)}}} X_{I_a} \given X_{\{a\} \cup J}.$$
    \item ``Having no confounding bias'' is invariant under interventional equivalence of iSCMs.
    \item Colloquially, confounding bias is also often referred to as ``spurious dependence''.
  \end{itemize}
\end{frame}

\begin{frame}{Confounding (hierarchy)}
\begin{Prp}[\ref{prp:graph_unconfounded_scm_rung3_unconfounded}]
  Let $M = \lt \Sin, \Send, \Srnd, \CX, P, f \rt$ be a simple iSCM with causal graph $G(M)$.
  Let $a, b \in V$ with $a \ne b$.
  If $a$ and $b$ are unconfounded according to $G(M)$ then $a$ and $b$ are rung-3 unconfounded according to $M$.
\end{Prp}

\begin{Prp}[\ref{prp:rung3_unconfounded_implies_no_confounding_bias}]
  Let $M = \lt \Sin, \Send, \Srnd, \CX, P, f \rt$ be a simple iSCM.
  Let $a,b \in V$ be distinct endogenous variables such that $b$ is not a rung-3 cause of $a$ according to $M$.
  If $a$ and $b$ are rung-3 unconfounded according to $M$, then $a$ and $b$ have no confounding bias according to $M$.
\end{Prp}
\end{frame}

\begin{frame}{Confounding (discussion)}
\begin{itemize}
  \item How to define ``$a$ and $b$ have confounding bias'' if $a$ and $b$ are part of a causal cycle is at present an open research problem.
  \item The assumption that two variables have no confounding bias is often (perhaps too often) made in practice.
  \item If one doesn't rule out possible confounding bias, one can still obtain bounds on causal predictions (see Theorem~\ref{thm:natural_bounds}).
\end{itemize}
\end{frame}
